% ECONOMICS_PROBLEM: {"id": "E62-552", "title": "State-dependent fiscal multipliers", "jel_code": "E62", "source_id": "OP-0014"}
% !TeX program = xelatex
\documentclass[11pt,a4paper]{article}
\usepackage[margin=23mm]{geometry}
\usepackage{fontspec}
\setmainfont{Latin Modern Roman}
\usepackage{amsmath,amssymb,mathtools}
\usepackage{xcolor,enumitem,booktabs,longtable,array,xurl,hyperref}
\definecolor{muted}{HTML}{53636C}
\hypersetup{colorlinks=true,urlcolor=blue,linkcolor=blue}
\setlength{\parindent}{0pt}
\setlength{\parskip}{5pt}
\newcommand{\R}{\mathbb R}
\newcommand{\E}{\mathbb E}
\newcommand{\N}{\mathbb N}
\newcommand{\ind}{\mathbf 1}
\newcommand{\Var}{\operatorname{Var}}
\newcommand{\Cov}{\operatorname{Cov}}
\newcommand{\supp}{\operatorname{supp}}
\DeclareMathOperator*{\argmin}{arg\,min}
\DeclareMathOperator*{\argmax}{arg\,max}
\newcommand{\entrymeta}[3]{{\sffamily\small\color{muted}\textbf{\texttt{#1}}\quad|\quad #2\par #3\par}}
\newcommand{\Problemref}[1]{\texttt{#1}}
\newcommand{\JELref}[2]{\texttt{#2} (JEL \texttt{#1})}
\begin{document}
\section*{State-dependent fiscal multipliers}
\textbf{Problem ID:} \texttt{E62-552}\quad \textbf{JEL:} \texttt{E62}\par
% BEGIN ECONOMICS STATEMENT
\label{problem:OP-0014}
\label{atlas:prob:A4}

\noindent\textbf{Original atlas ID:} A4. \textbf{Formulation:} editorial specification of a research family.

\subsection*{Definitions and assumptions}

Let $S_t$ be a prespecified state measured before a fiscal intervention, and let $\varepsilon_t=a$ denote a structurally defined spending disturbance. Fix the monetary rule, financing rule, openness assumptions, sample population, and real-price units. For horizons $h\ge0$, let $Y_{t+h}(a)$ and $G_{t+h}(a)$ be potential aggregate output and government purchases under that intervention, allowing other endogenous variables and subsequent shocks to follow the declared regime. Define $\Delta_Y(h,s,a)=\mathbb E[Y_{t+h}(a)-Y_{t+h}(0)\mid S_t=s]$, and analogously $\Delta_G$. For positive discount weights $d_h$ and finite $H$, the cumulative multiplier is
\[
\mathcal M_H(s,a)=\frac{\sum_{h=0}^Hd_h\Delta_Y(h,s,a)}
 {\sum_{h=0}^Hd_h\Delta_G(h,s,a)},
\]
provided its denominator is nonzero.

\subsection*{Formal research question}

Identify $\mathcal M_H(s,a)$, or a sharp identified set for it, across prespecified states describing slack, monetary accommodation, debt, and spending composition. Distinguish a finite-shock multiplier from the local derivative obtained by replacing $\Delta_Y$ and $\Delta_G$ with derivatives at $a=0$. Determine which instruments, narrative timing assumptions, or quasi-experiments identify the requisite intervention law and its state interactions. Estimate uncertainty for the ratio, including weak-denominator cases, and compare specifications on common data and units. The state must not be redefined using post-intervention outcomes. Treatment anticipation, policy spillovers, and the financing response belong to the causal estimand and identification assumptions.

\subsection*{Interpretation and scope}

A multiplier is not solely a property of ``fiscal policy'': it depends on the intervention path, horizon, conditioning state, financing, monetary reaction, and output measure. Spending news and current purchases are different treatments when agents anticipate future fiscal actions. A regional estimated response may omit aggregate monetary or tax adjustments, so it is not automatically a national multiplier. Linear interactions can approximate state dependence but do not establish invariance under a large emergency intervention. Defining states before the shock avoids conditioning on a variable caused by the treatment. When the denominator is weakly identified or may cross zero, a conventional narrow delta-method interval can be invalid; a set-valued ratio or appropriate weak-identification method is required.

\subsection*{Source audit and status}

Blanchard--Perotti (2002), Ramey (2011), and Auerbach--Gorodnichenko (2012) are verified through their primary journal pages. The first uses institutional timing restrictions, the second emphasizes anticipation and identification timing, and the third estimates state-dependent responses. They already provide important restricted answers, so the question cannot be treated as wholly unstudied. Their differences motivate a common-estimand comparison rather than a universally correct numerical multiplier. The potential-outcome definitions and ratio target here are editorial formalizations. Remaining uncertainty is an empirical and identification agenda for explicit regimes and interventions, not a single theorem certified open by the historical references.

\subsection*{References}

\begin{enumerate}

\item[S1] Olivier Blanchard and Roberto Perotti (2002). \emph{An Empirical Characterization of the Dynamic Effects of Changes in Government Spending and Taxes on Output}. Quarterly Journal of Economics 117(4), 1329--1368. \url{https://academic.oup.com/qje/article-abstract/117/4/1329/1875961}. \textit{Locator:} abstract; institutional timing identification. \textit{Establishes:} Fiscal responses conditional on institutional identifying assumptions. \textit{Does not establish:} state-invariant or universally causal multipliers.

\item[S2] Valerie A. Ramey (2011). \emph{Identifying Government Spending Shocks: It's All in the Timing}. Quarterly Journal of Economics 126(1), 1--50. \url{https://academic.oup.com/qje/article-abstract/126/1/1/1902509}. \textit{Locator:} abstract; narrative timing comparison. \textit{Establishes:} How anticipation and timing change estimated fiscal responses. \textit{Does not establish:} a single multiplier valid under every policy regime.

\item[S3] Alan J. Auerbach and Yuriy Gorodnichenko (2012). \emph{Measuring the Output Responses to Fiscal Policy}. American Economic Journal: Economic Policy 4(2), 1--27. \url{https://www.aeaweb.org/articles?id=10.1257/pol.4.2.1}. \textit{Locator:} abstract; state-dependent fiscal responses. \textit{Establishes:} Evidence on fiscal effects across economic conditions. \textit{Does not establish:} a restriction-free answer to every state-dependent multiplier question.

\end{enumerate}

\subsection*{Common conventions: Source mathematical conventions}

Unless an entry states otherwise, agent and firm sets are finite. State and action spaces are measurable, strategies and outcome maps are measurable, probability laws are specified, and expectations exist for the targets under discussion. Infinite-horizon discount factors lie in $(0,1)$ and discounted objectives are finite. Norms, tolerances and complexity input sizes must be specified for computational comparisons. Constants or primitives that appear locally are inputs, not empirical facts asserted by this catalogue.

\textbf{Identification.} A statistical model $m\in\mathcal M$ implies an observable law $P_m$ and target $\theta(m)$. At law $P$, its identified set is
\[
\Theta(P;\mathcal M)=\{\theta(m):m\in\mathcal M,\ P_m=P\}.
\]
Point identification holds when this set is a singleton, or equivalently when $P_{m_1}=P_{m_2}$ implies $\theta(m_1)=\theta(m_2)$ for all compatible models. Sharp bounds mean bounds attained, or approached when the boundary is not attained, by compatible models. Writing this set is a definition, not a solution: the substantive task is to establish informative restrictions and derive or estimate the set. An unrestricted-confounding model may give only trivial bounds.

\textbf{Inference.} Identification, estimability and valid uncertainty quantification are separate questions. Fix $\alpha\in(0,1)$. A confidence procedure $C_n$ over a declared law class $\mathcal P$ has asymptotic uniform coverage at level $1-\alpha$ when
\[
\liminf_{n\to\infty}\inf_{P\in\mathcal P}\Pr_P\{\theta(P)\in C_n\}\geq1-\alpha.
\]
For set-identified targets the coverage event must be defined separately, for example coverage of the full identified set. If an entry uses identification-indexed models instead of uniquely indexed laws, the same coverage convention is applied over those models. Pointwise limits do not establish uniform validity, and asymptotic validity does not guarantee finite-sample precision. Sample sizes, dependence structures and weak-identification sequences are part of each application.

\textbf{Counterfactuals.} $Y_i(a)$ is a potential outcome under a specified intervention, including its timing, implementation and versions. It is not $Y_i$ conditional on voluntarily choosing $a$. A full intervention vector permits interference; shorthand scalar treatment notation does not silently impose no interference. Assignment, selection, attrition, anticipation and measurement must be specified. A claim about an instrument's local effect is not automatically a population policy effect.

\textbf{Economic equilibrium.} A counterfactual model includes best responses, constraints, feasibility, market clearing, state transitions and expectations consistent with the stated information. When several equilibria exist, either a declared selector is an input or the target is set-valued. Comparisons across policy regimes must use the stated selection convention. Reduced-form relationships are not presumed invariant after an equilibrium-changing intervention.

\textbf{Optimization and welfare.} Optimization is over the stated feasible set. If attainment is not established, $\sup$ or $\inf$ and $\varepsilon$-optimal choices replace an asserted maximizer or minimizer. For example, with value $V=\sup_{a\in\mathcal A}W(a)<\infty$, an $\varepsilon$-optimal set is $\{a\in\mathcal A:W(a)\geq V-\varepsilon\}$ for $\varepsilon>0$. Benefits, costs, risk and health or environmental outcomes can be aggregated only using declared units and conversion weights. Interpersonal utility weights, the treatment of fiscal transfers and foreign profits, discounting, and the behavioral welfare criterion are explicit inputs. A comparative-static derivative additionally requires existence and appropriate differentiability of the selected counterfactual.

\textbf{Sources.} Local labels [S1], [S2], and so forth restart in every question. Locators identify the relevant abstract, model, section or result at the level actually checked; a bibliographic or abstract-level verification is not represented as inspection of an entire proof. Working papers are distinguished from later publications and changing versions. Source summaries are paraphrased. All URL checks and editorial formulations refer to the audit date stated above.
% END ECONOMICS STATEMENT
\end{document}
